% Generated draft: requires mathematical and source review.
\documentclass[11pt]{article}
\usepackage[margin=25mm]{geometry}
\usepackage{fontspec}
\setmainfont{Latin Modern Roman}
\usepackage{amsmath,amssymb,mathtools,unicode-math}
\setmathfont{Latin Modern Math}
\usepackage{enumitem,needspace,xurl,hyperref}
\setlength{\parindent}{0pt}
\setlength{\parskip}{5pt}
\setlength{\emergencystretch}{3em}
\newcommand{\N}{\mathbb{N}}
\newcommand{\Z}{\mathbb{Z}}
\newcommand{\Q}{\mathbb{Q}}
\newcommand{\R}{\mathbb{R}}
\newcommand{\C}{\mathbb{C}}
\newcommand{\F}{\mathbb{F}}
\newcommand{\A}{\mathbb{A}}
\newcommand{\PP}{\mathbb{P}}
\newcommand{\E}{\mathbb{E}}
\newcommand{\eps}{\varepsilon}
\newcommand{\dd}{\,\mathrm{d}}
\begin{document}
\section*{UnsolvedMath ID: 700642}
\textbf{Source problem code:} catalogue entry
\subsection{Definitions and mathematical statement}
Let \(G\) be a finite group. An involution is an element of order exactly \(2\). Write \(t^x=x^{-1}tx\), \(t^G=\{t^x:x\in G\}\), and let \(\langle S\rangle\) denote the subgroup generated by a subset \(S\subseteq G\). For a finite group \(H\), let \(O_2(H)\) be its largest normal \(2\)-subgroup and \(O_{2'}(H)\) its largest normal subgroup of odd order. Write \(C_G(a)=\{g\in G:ga=ag\}\), \(C_D(t)=D\cap C_G(t)\), and \(N_G(X)=\{g\in G:g^{-1}Xg=X\}\).

\begin{enumerate}
\item[(a)] Let \(t\in G\) be an involution and define
\[
D=t^G\cup\{t^xt^y:x,y\in G,\ |t^xt^y|=2\}.
\]
Suppose that
\[
D\cap O_2(G)=\varnothing
\quad\text{and}\quad
\forall d\in C_D(t),\qquad t\in O_2(C_G(d)).
\]
Every element of \(D\) is an involution, so the second condition is precisely the source's condition on every involution in \(C_D(t)\). Prove that \(D=t^G\). The source additionally states that the structure of \(\langle D\rangle\) is known in this case, but supplies neither that description nor a supporting citation.

\item[(b)] Independently of the hypotheses in (a), let \(t\) be an involution of a finite group \(G\). The source uses \(Z^*(H)\) without defining it. Under the usual finite-group convention, if \(\pi_H:H\to H/O_{2'}(H)\) is the quotient map, then
\[
Z^*(H)=\pi_H^{-1}\bigl(Z(H/O_{2'}(H))\bigr).
\]
The following formulation is conditional on this being the intended meaning of \(Z^*\). Suppose that, for every subgroup \(X\leq G\),
\[
\left(1<|X|,\quad |X|\text{ is odd},\quad t\in N_G(X),\quad
\exists x\in X:\ tx\ne xt\right)
\ \Longrightarrow\ 
t\in Z^*(N_G(X)).
\]
Determine the structure of \(\langle t^G\rangle\) for all such pairs \((G,t)\). The source describes this as a significantly more general question; it does not specify the desired form of the structural description.
\end{enumerate}
\subsection{Short English statement}
Part (a) asks whether the stated conditions on centralizers force every involution that is a product of two conjugates of \(t\) to be itself conjugate to \(t\). Part (b) asks for the structure of the subgroup generated by all conjugates of \(t\), under a condition involving every nontrivial odd-order subgroup that \(t\) normalizes but does not centralize. The meaning of \(Z^*\) requires confirmation.
\subsection{Sources}
\begin{itemize}
\item Local UnsolvedMath registry, permanent ID \(700642\), exact source problem code \texttt{catalogue entry}, title \emph{Kourovka Notebook Problem 10.27}; supplied upstream revision \texttt{d371ee917ac9647fbb06cd6fc529deec6f3dda82}.
\item Attributed in the supplied record to A.~A.~Makhn\"ev, proposed in 1986. Source attribution: \emph{The Kourovka Notebook (Unsolved Problems in Group Theory)}, 10th issue (1986), problem 10.27; 21st issue edition (2026), identified in the record as arXiv:1401.0300v48. The collection description names E.~I.~Khukhro and V.~D.~Mazurov as editors.
\item The registry labels the problem \emph{open}. This is a supplied status label, not an independently verified mathematical or literature assessment. No source was independently consulted, and no separate citation is supplied for the claim that the structure in part (a) is known.
\end{itemize}
\end{document}
